\section{Runtime Analysis}
\label{sec:runtime}

\subsection{Theoretical Scaling}
\label{sec:runtime:theory}

The \texttt{redist} pipeline's runtime is dominated by the METIS $k$-way
partitioning step. METIS's multilevel scheme runs in $O(n \log n)$ time
for a graph with $n$ vertices \citep{karypis1998fast}. Moving from tract
($n_T$) to block-group ($n_{BG}$) resolution multiplies the vertex count
by approximately:
\[
  r = \frac{n_{BG}}{n_T} = \frac{220{,}000}{73{,}000} \approx 3.0.
\]
The theoretical runtime ratio is $r \log(r \cdot n_T) / \log(n_T) \approx 3.0$
for large $n_T$, consistent with the observed $\approx 3$--$4\times$ factor.

For the block-to-block-group ratio:
\[
  r' = \frac{n_{\rm block}}{n_{BG}} = \frac{8{,}100{,}000}{220{,}000} \approx 36.8,
\]
giving a theoretical block-to-block-group runtime ratio of $\sim\!40\times$.
Block-level redistricting for a large state would take hours per state, making
it practical only for small states (NH, WY, VT) where $n_{\rm block}$ is small.

\subsection{Observed Runtimes by State}
\label{sec:runtime:observed}

Table~\ref{tab:runtime-comparison} shows observed runtimes for the three
focal states at both resolutions.

\begin{table}[h]
\centering\small
\begin{tabular}{lrrrrrr}
\toprule
State & $k$ & $n_T$ & $n_{BG}$ & Tract (s) & Block-group (s) & Ratio \\
\midrule
Washington & 98  & 1{,}458  & 4{,}874   & 45  & 180  & 4.0$\times$ \\
Texas      & 150 & 5{,}265  & 15{,}811  & 95  & 380  & 4.0$\times$ \\
California & 80  & 8{,}997  & 25{,}914  & 140 & 520  & 3.7$\times$ \\
\bottomrule
\end{tabular}
\caption{Observed runtimes: tract vs.\ block-group resolution (single seed,
  single CPU core). The block-group/tract runtime ratio is approximately
  3.7--4.0$\times$ for all three states, consistent with the $\sim\!3\times$
  vertex count ratio and $O(n \log n)$ METIS scaling.}
\label{tab:runtime-comparison}
\end{table}

\subsection{Total Pipeline Runtime: All 50 States}
\label{sec:runtime:total}

For the F.1 50-state house sweep with adaptive resolution selection
(block-group for 12 states, tract for 38 states), the total single-core runtime is
approximately 68 minutes (mean 82 seconds per state). With 4 parallel workers,
this reduces to approximately 18 minutes.

The congressional pipeline runtime for comparison (50 states, tract resolution)
is approximately 22 minutes single-core. The state legislative pipeline is
$\sim\!3\times$ slower due to the block-group resolution states and the
generally larger $k$ values.

\subsection{Scaling to Block Resolution}
\label{sec:runtime:block}

For the five states requiring block resolution (NH, WY, VT, SD, ND), estimated
block-resolution runtimes are:
\begin{itemize}
  \item New Hampshire: $\sim\!4{,}100$ seconds ($\sim\!68$ minutes) at block
    resolution ($n_{\rm block} \approx 36{,}000$).
  \item Wyoming: $\sim\!800$ seconds at block resolution
    ($n_{\rm block} \approx 12{,}500$).
\end{itemize}
These runtimes are feasible for a one-time run but prohibitive for multi-seed
sensitivity analysis. The F.1 paper uses block-group resolution for these states
as a practical compromise, noting the limitation.
